\chapter{C1: Partizipien und ihre Verwendung}
\label{cha:C1Partizipien}

\section{Partizip I (Partizip Präsens)}
\label{sec:C1PartizipI}

Bildung: \textbf{Infinitiv} + -d

\begin{center}
\begin{tabular}{|l|c|c|}
\hline
Infinitiv & Partizip I & Bedeutung \\ \hline
lesen & lesend & reading \\ \hline
arbeiten & arbeitend & working \\ \hline
laufen & laufend & running \\ \hline
\end{tabular}
\end{center}

Verwendung:
\begin{itemize}
    \item \textbf{als Adjektiv}: das lesende Buch (the book that is reading)
    \item \textbf{adverbial}: while doing something
    \item \textbf{substantiviert}: der Lesende (the reader)
\end{itemize}

\section{Partizip II als Attribut}
\label{sec:C1PartizipIIAttr}

\begin{center}
\begin{tabular}{|l|c|c|}
\hline
Aktiv & Passiv & \\ \hline
das lesende Buch & das gelesene Buch & the book being read \\ \hline
der arbeitende Student & der gearbeitete Student & the student working \\ \hline
\end{tabular}
\end{center}

\begin{tippbox}{Merke}
Partizip II als Attribut muss dekliniert werden!
\end{tippbox}

\section{Erweiterte Attribute}
\label{sec:C1Erweitert}

Partizipien können erweitert werden:

\begin{center}
\begin{tabular}{|l|c|}
\hline
Einfach & Erweitert \\ \hline
das gelesene Buch & das gestern gelesene Buch \\ \hline
der gelehrte Student & der an der Universität gelehrte Student \\ \hline
\end{tabular}
\end{center}

\chapterend

\chapter{C1: Konnektoren und Diskursmarker}
\label{cha:C1Konnektoren}

\section{Argumentation}
\label{sec:C1Argumentation}

\begin{center}
\begin{tabular}{|l|c|c|}
\hline
Konnektor & Funktion & Beispiel \\ \hline
zunächst / erstens & Reihenfolge & \textbf{Zunächst} möchte ich... \\ \hline
außerdem / darüber hinaus & Addition & \textbf{ Außerdem} muss man... \\ \hline
jedoch / aber & Gegensatz & \textbf{Jedoch} bin ich der Meinung... \\ \hline
folglich / demnach & Schluss & \textbf{Folglich} ergibt sich... \\ \hline
beispielsweise / etwa & Beispiel & \textbf{Beispiel} sei genannt... \\ \hline
\end{tabular}
\end{center}

\section{Textkohäsion}
\label{sec:C1Textkohaesion}

\begin{center}
\begin{tabular}{|l|c|c|}
\hline
Funktion & Ausdruck & Beispiel \\ \hline
Verweis auf Vorheriges & dies-, solch-, solcher- & \textbf{Dies} zeigt, dass... \\ \hline
Zusammenfassung & zusammenfassend, insgesamt & \textbf{Zusammenfassend} lässt sich sagen... \\ \hline
Neuigkeit & im Folgenden, im Folgenden & \textbf{Im Folgenden} werde ich... \\ \hline
\end{tabular}
\end{center}

\chapterend

\chapter{C1: Nominalisierung und Verbalisierung}
\label{cha:C1Nominalisierung}

\section{Nominalisierung (Verb $\rightarrow$ Nomen)}
\label{sec:C1Nominalisierung}

\begin{center}
\begin{tabular}{|l|c|c|}
\hline
Verb & Nomen & Beispiel \\ \hline
entscheiden & die Entscheidung & eine wichtige Entscheidung treffen \\ \hline
analysieren & die Analyse & eine Analyse durchführen \\ \hline
verbessern & die Verbesserung & eine Verbesserung erzielen \\ \hline
entwickeln & die Entwicklung & eine Entwicklung vorantreiben \\ \hline
lösen & die Lösung & eine Lösung finden \\ \hline
\end{tabular}
\end{center}

\begin{infobox}{Merke}
Nominalisierung macht den Text formeller und akademischer.
\end{infobox}

\section{Verbalisierung (Nomen $\rightarrow$ Verb)}
\label{sec:C1Verbalisierung}

\begin{center}
\begin{tabular}{|l|c|c|}
\hline
Nomen & Verb & Beispiel \\ \hline
die Analyse & analysieren & Wir analysieren die Daten. \\ \hline
die Entscheidung & entscheiden & Sie entscheidet sich. \\ \hline
die Entwicklung & entwickeln & Sie entwickeln ein Programm. \\ \hline
\end{tabular}
\end{center}

\chapterend